\begin{longtable}{@{}p{0.07\textwidth} p{0.07\textwidth} p{0.20\textwidth} p{0.11\textwidth} p{0.45\textwidth}@{}}
\caption{\paramscaption}\label{\paramslabel}\\
\toprule
\textbf{Symbol} & \textbf{Value} & \textbf{Meaning} & \textbf{Source} & \textbf{What fixes it, or what assuming it costs} \\
\midrule
\endfirsthead
\multicolumn{5}{@{}l}{\emph{Table \thetable, continued from the previous page}}\\
\toprule
\textbf{Symbol} & \textbf{Value} & \textbf{Meaning} & \textbf{Source} & \textbf{What fixes it, or what assuming it costs} \\
\midrule
\endhead
\midrule
\multicolumn{5}{r@{}}{\emph{continued on the next page}}\\
\endfoot
\bottomrule
\endlastfoot
\multicolumn{5}{@{}l}{\emph{The local identification model}} \\
\addlinespace[2pt]
$\bar m$ & 10 & operating point: the level management holds the patient at & \textsc{assumed} & illustrative, and every design quantity is reported relative to it, so its absolute value cancels out of the conclusions. On a real scale: the translational model puts the untreated plasma 3-HP baseline at $60.1$~$\mu$M and the fully corrected floor near $0.1$~$\mu$M~\cite{pkpd_translational}, so a managed patient sits between the two and $\bar m$ is read in $\mu$mol/L \\
$\theta$ & 3000 & residual capacity of the enzymatic channel & \textsc{assumed} & illustrative, and the quantity the paper says cannot be identified, so no conclusion may depend on its value. Residual PCC activity is published only as a fraction of control, under $5\%$ in PA, never as an absolute flux, so $\theta$ is an order-of-magnitude quantity and is confounded with $V_{\max}$ by construction. It is a \emph{capacity} and not a flux: at $\bar m/K_m = 0.1$ the channel runs at $9.1\%$ of $V_{\max}$, carrying $273$ against a load of $323$, so it is not comparable with a load figure directly \\
$K_m$ & 100 & half-saturation of the enzymatic channel, at the plasma level & \textsc{assumed} & illustrative, and \emph{not} the in vitro Michaelis constant of PCC ($0.29$ mM for propionyl-CoA~\cite{pcc_liver_kinetics}), a different molecule in a different compartment. The ratio $\bar m/K_m = 0.1$ is what the excitation condition is stated against \\
$k_e$ & 5 & alternative clearance rate constant & \textsc{assumed} & illustrative, chosen so the enzymatic channel carries most of the flux, which is the favourable case for identification. \textbf{This is the priority gap}: Section~\ref{sec:designsec} needs it within a factor of about three, and no direct measurement exists. The nearest route is the first-order 3-HP elimination implicit in the published indirect-response model~\cite{pkpd_translational} \\
$\mtol$ & 40 & tolerable metabolite level & \textsc{assumed} & a clinical convention, not a property of the patient. The width of the undetermined window runs $1.2\%$ to $67\%$ as it moves, so this is the most expensive assumption in the local model. No numeric plasma 3-HP target is published; the only quantitative anchor is that a $50\%$ lower 3-HP carries a decompensation rate ratio of $0.7$~\cite{natural_history} \\
$\spk/s_{\mathrm{base}}$ & 1.50 & illness amplitude, as a fold of the patient's own baseline load & \textsc{assumed} & stated as a fold rather than absolutely, because an absolute figure disagreed between sections. Model C's own fold sits at 1.18. Muscle proteolysis during illness bounds this from above: 3-methylhistidine excretion rises by $227\%$ to $325\%$ in sepsis and trauma, so $1.50$ is a conservative value for a moderate intercurrent illness \\
$\kappa$ & not assigned & gain from recorded carnitine dose onto $k_e$ & \textsc{assumed} & never assigned a value. Only the product $\kappa u_c$ enters, as a lumped increment on $k_e$, and every figure varies that product directly. The ill-conditioning result of supplement \S S1 is conditional on it: an unknown $\kappa$ leaves the carnitine route two equations in four unknowns. The \emph{direction} is established, carnitine supplementation raising propionylcarnitine excretion~\cite{carnitine_excretion_pa}, but no dose-response slope is published, so the magnitude cannot be assigned \\
$\lambda$ & $\ln 2/3.5$ d$^{-1}$ & decay rate of the therapeutic increment & \textsc{literature} & a $3.5$-day half-life, so the increment has decayed sixteen-fold by the next dose and the inter-dose recovery is an excursion rather than a ripple. \textbf{Corroborated}: the translational PK/PD model estimates a PCC-protein degradation rate of $0.0075$~h$^{-1}$~\cite{pkpd_translational}, a half-life of $3.85$~d, so this is within $10\%$ of a published value rather than free. It was previously named and never assigned, which made Section~\ref{sec:collapse} unrunnable; its only role is to set how much excursion a fixed schedule creates, and Table~\ref{tab:excursion} reports the answer \\
$T_{\mathrm{dose}}$ & 14 d & interval between therapeutic doses & \textsc{literature} & the biweekly schedule selected for mRNA-3927 by model-informed dose selection~\cite{dose_selection_poster}; the phase 1/2 trial dosed every two or three weeks~\cite{trial_phase12}. Adopted rather than assumed \\
$D$ & 6000 & size of one therapeutic increment & \textsc{assumed} & chosen so the marker falls by about half at peak effect, the order of the reductions reported in the phase 1/2 interim analysis~\cite{trial_phase12}. Only the excursion it produces is used. It is recoverable in principle: the published dose-to-marker transfer function~\cite{pkpd_translational} turns a measured per-dose fall in 3-HP into a delivered increment \\
$\sigma$ & a family & standard deviation of the ordinate of the fit & \textsc{assumed} & deliberately not assigned. Every design conclusion is reported as a function of it, because nothing in this paper bounds how well the substrate load is recorded. The assay contributes $b\,\sigma_m$ to the same total \\
$\mathrm{CV}$ & 5\% & analytical variability of the marker assay & \textsc{assumed} & illustrative, and a \emph{percentage}: it is never absolute flux. It enters the ordinate as $b\,\mathrm{CV}\,\bar m$. This is the \emph{analytical} figure: the breath test reports repeatability of $6.9\%$ in controls and $12.8\%$ in patients~\cite{breath_test_mma}. Biological variability is far larger, the translational model carrying $51\%$ between subjects on the 3-HP baseline~\cite{pkpd_translational} \\
\addlinespace
\midrule
\multicolumn{5}{@{}l}{\emph{The parent model and its equilibrium reduction}} \\
\addlinespace[2pt]
$\theta$ & 300 & residual capacity, parent model & \textsc{assumed} & illustrative, about $10\%$ of normal, chosen to sit above $\thc^{\,\mathrm{int}}$ so the patient has a margin to measure. The published figure is a fraction rather than a flux: residual PCC activity is under $5\%$ of control in PA, and converting that to an absolute $V_{\max}$ has never been done, which would need a PA-specific $1$-$^{13}$C-propionate breath test~\cite{breath_test_mma} \\
$K_m$ & 100 & half-saturation of the enzymatic channel & \textsc{assumed} & as above, and shared with the local model so the two can be compared \\
$K_c$ & 100 & half-saturation of the conjugation step & \textsc{assumed} & no measurement, and no reason a conjugation reaction should share an enzyme's half-saturation. Cost bounded in supplement \S S1. The nearest measured analogue is the renal carnitine transporter, $K_m$ $0.08$ mM external and $1.2$ mM internal~\cite{octn2_renal_transport}, which is a different step; no half-saturation for the conjugation itself is published \\
$M_i$ & 300 & metabolite level at which acylcarnitine competition doubles the non-conjugative loss & \textsc{assumed} & no measurement. Previously not written down anywhere, which is why the margin figure could not be reproduced from the paper. Cost bounded in supplement \S S1. The mechanism is measured, acylcarnitines being the most effective inhibitors of that transporter~\cite{octn2_renal_transport}, but no propionylcarnitine inhibition constant is published, so the level at which competition doubles the loss is not available as a number \\
$\rho_0$ & 0.50 & basal non-conjugative loss relative to conjugation, $\delta_0/k$ & \textsc{assumed} & no direct measurement. Bounded above by $1$, since the carnitine drain in this disease is dominated by propionylcarnitine, and below by the measured direction of free-carnitine excretion: in PA and MMA the urinary excretion of free carnitine falls while short-chain acylcarnitine excretion rises~\cite{carnitine_excretion_pa} \\
$R_0$ & 1122.0 & carnitine replenishment & \textsc{derived} & \textbf{not free}. Fixed by requiring the alternative route alone to peak at $500$ against a baseline load of $600$, so that a patient with $\theta = 0$ cannot be maintained. Changing $K_c$, $M_i$ or $\rho_0$ moves that peak and $R_0$ must be re-derived \\
$s_{\mathrm{base}}$ & 600 & baseline propiogenic load & \textsc{assumed} & illustrative, and sets the scale everything else in the parent model is quoted against. Whole-body propionate turnover is published at roughly $250$ to $400$~$\mu$mol/kg/day, partitioned as about half from amino-acid catabolism, a fifth from gut fermentation and the rest from odd-chain fatty acids; a protein-restricted diet lowers the first two deliberately. The value here is not in those units and is not claimed to be \\
$k$ & cancels & conjugation rate constant & \textsc{derived} & drops out of the quasi-steady reduction of supplement \S S1 exactly, so the reduced clearance does not depend on it and it is a free scale in the two-state system \\
\end{longtable}